\section{总结}\label{sec:summary}

本工作中，我们用赝PDF方法首次从格点QCD计算了$\pi$介子胶子PDF，并研究了其$\pi$介子质量与晶格间距依赖。
我们在具有$N_f = 2+1+1$味高度改进 staggered 夸克（HISQ）的系综上使用 clover 价费米子，涵盖两种晶格间距（$a\approx 0.12$与0.15~fm）和三种$\pi$介子质量（220、310与690~MeV）。
这些系综使我们得以考察$\pi$介子胶子PDF对$\pi$介子质量与晶格间距的依赖；
在这两种情形下，依赖与当前统计不确定度相比都显得较弱。

我们研究了重构拟合所用函数形式相关的系统误差，以及匹配中忽略夸克贡献造成的系统误差。
通过采用其他 PDF 工作中常用或提出的多种形式，考察了所设胶子PDF拟合形式的影响。
我们观察到把$xg(x, \mu)/\langle x \rangle_g$的拟合从单参数改为双参数形式影响较大，但在三参数处收敛。
这意味着双参数拟合对本计算已经足够，最终的$\pi$介子胶子PDF结果以双参数拟合结果给出。
我们用 xFitter 的$\pi$介子夸克PDF估计了夸克对$\pi$介子胶子RITD的贡献，
发现其贡献的系统误差小于统计误差的$10\%$。

如最终比较图所示，最轻$\pi$介子质量下的$\pi$介子胶子PDF在不确定度内于$x>0.2$与 JAM、于$x>0.5$与 xFitter 一致。
我们还以$(1-x)^C$的形式研究了$\pi$介子胶子PDF在大-$x$区的渐近行为：
在两种晶格间距、三种$\pi$介子质量下，我们的研究暗示$C>3$。
未来精度更高、系统误差控制更好的格点QCD $\pi$介子胶子PDF研究，与预期实验~\cite{Denisov:2018unj,Arrington:2021biu,Anderle:2021wcy,Denisov:2018unj}的结果一道纳入全局拟合分析时，将为$\pi$介子内的胶子成分提供最佳测定。

